Kombinovaný rozptylový diagram zachycující zároveň vztah pokrytí kolektivními smlouvami ke čtyřem výstupním proměnným (čistý příjem v~\ac{PPS}, Giniho koeficient, počet odpracovaných hodin, zaměstnanecký podíl vs.~\ac{OSVČ}) umožňuje simultánní evaluaci omnibusové hypotézy práce: silnější \ac{KV} instituce jsou systematicky asociovány s~lepšími výsledky napříč všemi sledovanými dimenzemi.
Každý z~dílčích grafů v~kombinaci potvrzuje parciální asociaci, ale kombinovaný pohled ukazuje, že jde o~koherentní systémový vzorec, nikoliv o~náhodné koincidence: země s~vysokým pokrytím \ac{KS} (AT, DK, BE, FR) konzistentně vykazují vyšší příjmy v~\ac{PPS}, nižší Gini, méně hodin práce za srovnatelný příjem a~nižší podíl \ac{OSVČ}.
ČR se ve všech čtyřech dílčích grafech umísťuje v~kvadrantu „nízké pokrytí / horší výsledky“, ačkoliv v~absolutních hodnotách není na úplném dně EU27.
Tento polysystémový argument je silnější než dílčí korelace: holistické zlepšení institucionálního rámce \ac{KV} (erga omnes + \ac{APZ}) by s~vysokou pravděpodobností zlepšilo celou sadu výstupů současně, protože sdílejí společný institucionální determinant (srov.~obr.~\ref{fig:coverage_income_scatter}, \ref{fig:coverage_gini_scatter}).

Kombinovaný rozptylový diagram zachycující zároveň vztah pokrytí kolektivními smlouvami ke čtyřem výstupním proměnným (čistý příjem v~\ac{PPS}, Giniho koeficient, počet odpracovaných hodin, zaměstnanecký podíl vs.~\ac{OSVČ}) umožňuje simultánní evaluaci omnibusové hypotézy práce: silnější \ac{KV} instituce jsou systematicky asociovány s~lepšími výsledky napříč všemi sledovanými dimenzemi.
Každý z~dílčích grafů v~kombinaci potvrzuje parciální asociaci, ale kombinovaný pohled ukazuje, že jde o~koherentní systémový vzorec, nikoliv o~náhodné koincidence: země s~vysokým pokrytím \ac{KS} (AT, DK, BE, FR) konzistentně vykazují vyšší příjmy v~\ac{PPS}, nižší Gini, méně hodin práce za srovnatelný příjem a~nižší podíl \ac{OSVČ}.
ČR se ve všech čtyřech dílčích grafech umísťuje v~kvadrantu „nízké pokrytí / horší výsledky“, ačkoliv v~absolutních hodnotách není na úplném dně EU27.
Tento polysystémový argument je silnější než dílčí korelace: holistické zlepšení institucionálního rámce \ac{KV} (erga omnes + \ac{APZ}) by s~vysokou pravděpodobností zlepšilo celou sadu výstupů současně, protože sdílejí společný institucionální determinant (srov.~obr.~\ref{fig:coverage_income_scatter}, \ref{fig:coverage_gini_scatter}).

Kombinovaný rozptylový diagram zachycující zároveň vztah pokrytí kolektivními smlouvami ke čtyřem výstupním proměnným (čistý příjem v~\ac{PPS}, Giniho koeficient, počet odpracovaných hodin, zaměstnanecký podíl vs.~\ac{OSVČ}) umožňuje simultánní evaluaci omnibusové hypotézy práce: silnější \ac{KV} instituce jsou systematicky asociovány s~lepšími výsledky napříč všemi sledovanými dimenzemi.
Každý z~dílčích grafů v~kombinaci potvrzuje parciální asociaci, ale kombinovaný pohled ukazuje, že jde o~koherentní systémový vzorec, nikoliv o~náhodné koincidence: země s~vysokým pokrytím \ac{KS} (AT, DK, BE, FR) konzistentně vykazují vyšší příjmy v~\ac{PPS}, nižší Gini, méně hodin práce za srovnatelný příjem a~nižší podíl \ac{OSVČ}.
ČR se ve všech čtyřech dílčích grafech umísťuje v~kvadrantu „nízké pokrytí / horší výsledky“, ačkoliv v~absolutních hodnotách není na úplném dně EU27.
Tento polysystémový argument je silnější než dílčí korelace: holistické zlepšení institucionálního rámce \ac{KV} (erga omnes + \ac{APZ}) by s~vysokou pravděpodobností zlepšilo celou sadu výstupů současně, protože sdílejí společný institucionální determinant (srov.~obr.~\ref{fig:coverage_income_scatter}, \ref{fig:coverage_gini_scatter}).

Kombinovaný rozptylový diagram zachycující zároveň vztah pokrytí kolektivními smlouvami ke čtyřem výstupním proměnným (čistý příjem v~\ac{PPS}, Giniho koeficient, počet odpracovaných hodin, zaměstnanecký podíl vs.~\ac{OSVČ}) umožňuje simultánní evaluaci omnibusové hypotézy práce: silnější \ac{KV} instituce jsou systematicky asociovány s~lepšími výsledky napříč všemi sledovanými dimenzemi.
Každý z~dílčích grafů v~kombinaci potvrzuje parciální asociaci, ale kombinovaný pohled ukazuje, že jde o~koherentní systémový vzorec, nikoliv o~náhodné koincidence: země s~vysokým pokrytím \ac{KS} (AT, DK, BE, FR) konzistentně vykazují vyšší příjmy v~\ac{PPS}, nižší Gini, méně hodin práce za srovnatelný příjem a~nižší podíl \ac{OSVČ}.
ČR se ve všech čtyřech dílčích grafech umísťuje v~kvadrantu „nízké pokrytí / horší výsledky“, ačkoliv v~absolutních hodnotách není na úplném dně EU27.
Tento polysystémový argument je silnější než dílčí korelace: holistické zlepšení institucionálního rámce \ac{KV} (erga omnes + \ac{APZ}) by s~vysokou pravděpodobností zlepšilo celou sadu výstupů současně, protože sdílejí společný institucionální determinant (srov.~obr.~\ref{fig:coverage_income_scatter}, \ref{fig:coverage_gini_scatter}).

\input{texparts/python/scatter_combined}



